% Copyright 2026 auto-research contributors. MIT; see repository LICENSE.
\section{证据比较与竞争解释}
\subsection{判断需要哪些支持}
\textbf{示例判断：}读者应能从结论直接找到支持它的来源。
这是本样例的设计目标，尚未进行用户测试。
正式报告需要说明读过哪些原文、比较条件是否一致，以及证据能支持到哪里。

\begin{table}[htbp]
\centering
\caption{示例证据比较；内容为写作提示}
\begin{tabular}{p{.20\linewidth}p{.32\linewidth}p{.32\linewidth}}
\toprule
比较对象 & 可支持的判断 & 尚未排除的因素\\
\midrule
原始材料 & 指明原文位置及阅读深度 & 版本变化与遗漏附件\\
比较条件 & 说明控制项与差异 & 资源预算与选择偏差\\
替代解释 & 给出可辨别的预测 & 尚未执行的检验\\
\bottomrule
\end{tabular}
\end{table}

\subsection{局限与下一步}
本文件只检查中文、公式、表格、引用与链接的呈现。
它不证明模板适合全部报告长度，也不表示用户已经选定默认模板。
如把真实研究内容迁入模板，须保留原来的主张版本、证据边界和人工审核状态。

\section{具体待审核项}
\begin{enumerate}
\item 中文正文、标题和数学符号是否清楚，长时间阅读是否舒适？
\item 证据比较表是否便于核对支持范围与未排除因素？
\item 章节层次、目录及超链接是否符合预期的报告长度？
\item 是否选择此版式用于下一份报告？此项尚未决定。
\end{enumerate}
